\chapter{Statistische Mechanik und Thermodynamik}

    \section{Fragestellung}

\begin{center}
    \includegraphics[width=0.6\textwidth]{Abb/1_1.pdf}
\end{center}
Obwohl durch mikroskopische Theorien, wie der Quantenmechanik, die Beschreibung
von Systeme exakt gelingt, ist diese Methode nur für wenige Teilchen sinnvoll.
Einen praxistauglichen Ansatz liefert die statistische Mechanik. Hierbei wird
von einer mikroskopischen Beschreibung, links im Bild dargestellt durch einzelne
Teilchen im Topf, zu einer Makroskopischen, rechts durch große Untersysteme, die
selbst einige Millionen Teilchen beinhalten, übergegangen. Der Schritt hin zu
makroskopischen Messgrößen, die diese Untersysteme charaktierisieren, soll nun
der erste Schritt sein.

    \subsection{Viele mikroskopische Freiheitsgrade - Mikrozustände}

klassische Mechanik: $\{\vec{r}_i, \vec{p}_i\} \quad i = 1, \dots , N$\\
$\rightarrow$ 6N Freiheitsgrade

\begin{align*}
    &\dot{\vec{r}_j} = \frac{\partial H}{\partial \vec{p}_j} \left( \{
                      \vec{r}_i, \vec{p}_i \} \right) \quad j = 1, \dots, N\\
    &\dot{\vec{p}_j} = - \frac{\partial H}{\partial \vec{r}_j} \left( \{
                         \vec{r}_i, \vec{p}_i\}\right)
\end{align*}
